\section{Proceso de Cinemática} \label{sec:proceso_cinematica}

Aquí explicarán su código. Si quieren mostrar una parte, pueden hacerlo de la siguiente forma

% Si no quieres ponerle título al código, puedes dejarlo en blanco.
\begin{lstlisting}[language=Matlab, caption={}]
	function [q_sol, p_sol] = cinematica_inv(r, p_des, tol, max_iter, alpha, numMuestras)
\end{lstlisting}

Pero solo háganlo en partes muy específicas (las que van a explicar en ese momento). No copien todo el código ya que eso está en GitHub.

Si les sale el error \texttt{latexminted no se reconoce como un comando interno o externo, programa o archivo por lotes ejecutable}, deben tener instalado python y usar el siguiente comando.
\begin{lstlisting}[language=bash, caption={bash: Instalar minted en python con pip}]
	pip install latexminted==0.5.1
\end{lstlisting}

\subsection{Cinemática Directa}
Primero, se crea una matriz homogenea A0, la cual representará la posición inicial. 

\begin{lstlisting}[language=Matlab, caption={}]
	R_inicial = euler2rotMat(orientacion_inicial, secuencia);
	Matriz_cero = [0 0 0]
	A0 = [R_inicial posicion_inicial
	Matriz_cero 1]
\end{lstlisting}

Después, utilizando la tabla que contiene los parámetros de Denavit Hartenberg y los límites de nuestro robot, establecemos el robot y la trayectoria a seguir.

\begin{lstlisting}[language=Matlab, caption={}]
	dh = readtable('datos\tabla_DH\robotFinal.csv');
	robot = crear_robot(dh,A0);
periodo = 2; % Periodo

	[q,dq,ddq] = trayectoria_q(robot,t,periodo);
\end{lstlisting}

En este caso, las articulaciones del robot realizarán un movimiento repetitivo cada 2 segundos, ya que este es el valor definido en el periodo del movimiento cíclico.

Para observar el resultado de la animación y la creación del robot, ir a la \autoref{fig:pruebaDH_robot}

\subsection{Cinemática Diferencial}
Explicar las partes importantes del código de la cinemática diferencial.
Para ver los resultados, ir al \autoref{chap:resultados}: Resultados, o determinada figura.

\subsection{Cinemática Inversa}
Explicar las partes importantes del código de la cinemática inversa.
Para ver los resultados, ir al \autoref{chap:resultados}: Resultados, o determinada figura.

